% ═══════════════════════════════════════════════════════════════
% QR-CODE: Lernpfad Prueba Oral
% Klasse 10 Spätbeginner
% ═══════════════════════════════════════════════════════════════

\documentclass[11pt, a4paper]{article}

\usepackage[utf8]{inputenc}
\usepackage[T1]{fontenc}
\usepackage[ngerman]{babel}

\usepackage{helvet}
\renewcommand{\familydefault}{\sfdefault}

\usepackage[margin=2.5cm]{geometry}
\usepackage{qrcode}
\usepackage{hyperref}
\usepackage{xcolor}

\definecolor{spanischrot}{HTML}{C41E3A}

\hypersetup{
    colorlinks=true,
    linkcolor=spanischrot,
    urlcolor=spanischrot
}

\setlength{\parindent}{0pt}
\pagestyle{empty}

\begin{document}

\begin{center}

\vspace*{2cm}

{\LARGE\bfseries Lernpfad}\\[8pt]
{\Large Prueba oral -- Repaso}\\[4pt]
{\normalsize Klasse 10 · Spanisch · Niveau A1}

\vspace{2cm}

% QR-Code
\qrcode[height=6cm]{https://devbox42.github.io/sciencesim/spanisch/kl10/LP-REPASO-PRUEBA.html}

\vspace{1.5cm}

% Clickable Link
{\large\textbf{Link zum Lernpfad:}}\\[8pt]
\url{https://devbox42.github.io/sciencesim/spanisch/kl10/LP-REPASO-PRUEBA.html}

\vspace{2cm}

\textcolor{spanischrot}{\rule{8cm}{1pt}}

\vspace{1cm}

{\small Scanne den QR-Code mit deinem Smartphone\\oder klicke auf den Link oben.}

\vspace{1cm}

{\footnotesize\textit{¡Buena suerte con tu preparación!}}

\end{center}

\end{document}
